\documentclass[11pt]{article}
\usepackage[margin=1in]{geometry}
\usepackage[utf8]{inputenc}
\usepackage[T1]{fontenc}
\usepackage{lmodern}
\usepackage{amsmath,amssymb,amsthm}
\DeclareUnicodeCharacter{2200}{\ensuremath{\forall}}
\DeclareUnicodeCharacter{2203}{\ensuremath{\exists}}
\DeclareUnicodeCharacter{2192}{\ensuremath{\to}}
\DeclareUnicodeCharacter{2264}{\ensuremath{\le}}
\DeclareUnicodeCharacter{2260}{\ensuremath{\ne}}
\DeclareUnicodeCharacter{2227}{\ensuremath{\wedge}}
\DeclareUnicodeCharacter{00AC}{\ensuremath{\neg}}
\DeclareUnicodeCharacter{2115}{\ensuremath{\mathbb{N}}}
\usepackage[title]{appendix}
\usepackage[colorlinks=true,linkcolor=blue,citecolor=blue,urlcolor=blue,pdfauthor={P. Graham},pdftitle={Monotone classes of bounded game chromatic number have bounded arboricity}]{hyperref}
\newtheorem{theorem}{Theorem}[section]
\newtheorem{corollary}[theorem]{Corollary}
\newtheorem{lemma}[theorem]{Lemma}
\newtheorem{question}[theorem]{Question}
\theoremstyle{remark}
\newtheorem{remark}[theorem]{Remark}
\setlength{\emergencystretch}{3em}
\title{Monotone classes of bounded game chromatic number have bounded arboricity}
\author{P. Graham\\ \small\url{https://orcid.org/0009-0007-8493-1505}}
\date{November 19, 2026}
\begin{document}
\maketitle
\begin{abstract}
In the coloring game on a finite graph, Alice and Bob alternately color
uncolored vertices with one of $k$ colors so that adjacent vertices get
different colors. Alice moves first and wins if every vertex gets colored;
Bob wins if some uncolored vertex sees all $k$ colors among its neighbors.
The game chromatic number of a graph is the least $k$ for which Alice can
force a win. We prove that for every $K$ there is $d$ such that every graph
of average degree at least $d$ has a subgraph on which Bob wins the
$k$-color game for every $k\le K$. Consequently, if a class of graphs is
closed under taking subgraphs and its game chromatic numbers are bounded,
then the edges of each of its graphs can be partitioned into a bounded
number of forests, that is, the class has bounded arboricity. This answers
Zhu's Question~3 (1999) in the affirmative. The proof combines the theorem of Kühn and Osthus
that every graph of large enough average degree contains a subgraph of
large average degree with no 4-cycle, with a strategy for Bob that adapts
the weighting argument behind Beck's criterion for positional games: Bob
uses weights to keep some uncolored vertex that can still be blocked from
every color. The rest of the proof is checked in Lean, with the theorem of
Kühn and Osthus as an explicit hypothesis.

This work has not yet been reviewed by experts.
The proofs, the Lean formalization and the text were produced by AI models
(from OpenAI and Anthropic), working under the author's direction.
\end{abstract}
\medskip
\noindent\textit{Keywords:} graph coloring game, game chromatic number,
arboricity, monotone graph class, 4-cycle-free subgraph, Maker--Breaker
game.

\noindent\textit{MSC 2020:} 05C15, 05C57, 91A43.

\section{Introduction}

\subsection{Main results}
Graphs are finite and simple. In the \emph{$k$-color game} on a graph $G$,
Alice and Bob alternately choose an uncolored vertex and give it one of the
colors $1,\dots,k$ that no neighbor of the vertex has; Alice moves first.
The game ends when the player to move has no legal move, and Alice wins if
every vertex is then colored; otherwise Bob wins. We use the rules as
stated by Zhu \cite[pp.~1--2]{Zhu}; the game goes back to Brams, was
published by Gardner and was introduced for general graphs by Bodlaender
\cite[Section~1]{BGKZ}. The \emph{game chromatic number} $\chi_g(G)$ is the
least $k$ for which Alice has a winning strategy. A class of graphs is
\emph{monotone} if it is closed under taking subgraphs, that is, any vertex
subset with any subset of the edges among it. The
\emph{arboricity} of $G$ is the least number of forests into which its edge
set can be partitioned. The average degree of a graph with at least one
vertex is $2|E(G)|/|V(G)|$.

\begin{theorem}\label{thm:main}
For every integer $K\ge0$ there is an integer $d=d(K)$ such that every
graph with at least one vertex and average degree at least $d$ has a
subgraph $H$ on which Bob has a winning strategy in the $k$-color game for
every integer $k$ with $0\le k\le K$; equivalently, $\chi_g(H)>K$.
\end{theorem}

Whether a win for Alice with $k$ colors implies one with $k+1$ colors is
Zhu's Question~1 \cite[p.~11]{Zhu}, still open \cite[Section~8]{BGKZ},
\cite[Question~1.1]{Ho}; so a win for Bob with $K$ colors alone would leave
the smaller palettes to check.

\begin{corollary}\label{cor:arboricity}
For every integer $K\ge0$ there is an integer $d(K)$ with the following
property. If a class of graphs is closed under taking subgraphs and each of
its members has game chromatic number at most $K$, then each of its
members has arboricity at most $d(K)$.
\end{corollary}

This answers in the affirmative Zhu's Question~3, whose hereditary classes
are those ``closed under taking (not necessarily induced) subgraphs''
\cite[p.~12]{Zhu}:

\begin{question}[{Zhu \cite[Question 3]{Zhu}}]
``Suppose a hereditary class of graphs has bounded game chromatic number.
Is it true that the arboricity of this class of graphs is also bounded?''
\end{question}

Both results rest on a theorem of Kühn and Osthus \cite[Theorem~1]{KO}, in
the form of Theorem~\ref{thm:ko}; given it as a hypothesis, both are proved
in Lean (Section~\ref{sec:lean}).

\subsection{Proof overview}
Given a graph of very large average degree, the theorem of Kühn and Osthus
gives a subgraph $G_0$ of large average degree with no 4-cycle, and inside
it a nonempty vertex set $S$ each of whose vertices has many neighbors in
$S$ (Lemma~\ref{lem:mindeg}). Labeling $S$ at random and averaging
(Lemma~\ref{lem:extract}) gives disjoint sets $U_1,\dots,U_K\subseteq S$
and a set $A\subseteq S$ outside them, more than $2^K+1$ times as large as
$U=U_1\cup\dots\cup U_K$, such that every vertex of $A$ has a neighbor in
every $U_i$. Keeping one such edge from each $a\in A$ to each $U_i$ gives a
subgraph $H$ in which every $a\in A$ has exactly one neighbor $a_i$ in each
$U_i$. Two vertices of $A$ share at most one of these neighbors, or there
would be a 4-cycle.

Fix $k\le K$ and play the $k$-color game on $H$ with $U_1,\dots,U_k$; call
$i$ the \emph{designated} color of $U_i$. An uncolored $a\in A$ whose
$a_1,\dots,a_k$ carry colors $1,\dots,k$ can never be colored, since colors
are never removed, and Bob wins. Alice protects $a$ by making some $a_i$
\emph{dead}: coloring it with another color, or giving color $i$ to a
neighbor of it. A move changes the status of at most one vertex of $U$, but
may cost every target adjacent to it. While some vertex of $U$ is
\emph{free} (uncolored, with no neighbor of its designated color), Bob
gives one its designated color; afterwards he plays any legal move. He
maximizes a potential, the sum of $2^{t(a)}$ over the $a$ with no dead
$a_i$, where $t(a)$ counts the $a_i$ carrying their designated color,
placed by either player. As for the Erdős--Selfridge theorem
\cite[p.~299]{ES} and Beck's criterion \cite{Be}, Bob's greedy move gains
at least what Alice's next move can lose, up to an error the overlap bound
controls. The potential stays above a reserve proportional to the number
of free vertices, so once none is free, some uncolored $a$ sees all $k$
colors. Finally, if all subgraphs of $G$ have average degree below $d$, removing
vertices with fewer than $d$ remaining neighbors one at a time orders the
vertices so that each has fewer than $d$ earlier neighbors; giving these
edges distinct colors splits the edges into $d$ forests.

\subsection{Prior work and related questions}
Zhu's questions concern hereditary classes because a subgraph can have
larger game chromatic number than the graph itself \cite[p.~11]{Zhu}, as
Tuza and Zhu also note in their survey \cite[Section~1, p.~306]{TZ}.

Zhu's Question~2 asks the same with the game coloring number in place of
the arboricity. In Bodlaender's marking game, where Alice and Bob
alternately mark vertices, Alice tries to keep small the largest number of
previously marked neighbors of any vertex; the least bound she can
guarantee is the \emph{Go number}, and the game coloring number is one more
\cite[Section~2]{KT}. A yes answer to Question~2 would imply one to
Question~3 \cite[p.~12]{Zhu}. In 2001 Kierstead and Trotter gave ``an
example which shows that the game chromatic number of a class of graphs
may be bounded even when the Go number is not'' \cite[Section~6]{KT}, and
with a theorem of Dinski and Zhu
\cite[Theorem~2]{DZ} their Example~6.1 gives a monotone class with game
chromatic number at most 12 and unbounded Go number, a negative answer to
Question~2 (Appendix~\ref{app:related}). Zhu's Question~4 repeats
Conjecture~1 of Dinski and Zhu \cite{DZ}, that every monotone class of bounded game
chromatic number has bounded acyclic chromatic number, the fewest colors
in a proper coloring with no two-colored cycle. It remains open, and would
imply Corollary~\ref{cor:arboricity}, since bounded acyclic chromatic
number implies bounded arboricity \cite[p.~12]{Zhu}; Appendix~\ref{app:related}
discusses Dvořák's reduction of it.

The new step is Theorem~\ref{thm:potential}, which adapts the weighting
argument behind Beck's criterion for Maker--Breaker games \cite{Be},
\cite[Theorem~4.3]{Kr}. Unlike Maker, Bob must move legally, and a new
estimate bounds the weight Alice's next move can destroy, since two
targets share at most one vertex of $U$.


\section{Bob's potential strategy}\label{sec:potential}

\begin{theorem}\label{thm:potential}
Let $k\ge0$ and $D\ge0$ be integers, and let $G$ be a graph with a set $A$
of vertices and pairwise disjoint sets $U_1,\dots,U_k$ of vertices,
disjoint from $A$. Write $U=U_1\cup\dots\cup U_k$ and assume:
\begin{enumerate}
\item every neighbor of every vertex of $U$ lies in $A$;
\item every $a\in A$ has exactly one neighbor in $U_i$, written $a_i$, for
each $i$;
\item two distinct vertices of $A$ have at most one common neighbor in $U$;
\item every vertex of $U$ has at most $D$ neighbors in $A$, and
$2^k|U|+8D<8|A|$.
\end{enumerate}
Then Bob has a winning strategy in the $k$-color game on $G$.
\end{theorem}

The graph $G$ may have further vertices, in neither $A$ nor $U$, and
further edges not meeting $U$.

\begin{proof}
A \emph{position} is a partial proper coloring reached by legal moves.
In a position $c$, call a vertex $w\in U_i$ \emph{taken} if $c(w)=i$,
\emph{free} if $w$ is uncolored and no neighbor of $w$ has color $i$, and
\emph{dead} otherwise. Colors never change, so a taken vertex stays taken
and a dead vertex stays dead. Call $a\in A$ \emph{alive} if none of
$a_1,\dots,a_k$ is dead, and let $t(a)$ be the number of taken vertices
among them. Put
\[
\Phi(c)=\sum_{a\text{ alive}}2^{t(a)},\qquad
\Phi_y(c)=\sum_{a\text{ alive},\ a\sim y}2^{t(a)}\quad(y\in U),
\]
and let $F(c)$ be the number of free vertices.

(a) \emph{An alive vertex is uncolored.} If $a\in A$ has color $j$, then
$a_j$ is not taken (the coloring is proper) and not free (its neighbor $a$
has color $j$), so $a_j$ is dead.

(b) \emph{One move changes the status of at most one vertex of $U$.} Let a
move give the vertex $v$ the color $j$. If $v\in U_i$, then by (1) no
vertex of $U$ is adjacent to $v$, so only $v$ can change status: if it
was free, it becomes taken if $j=i$ and dead otherwise, and if it was dead,
it stays dead. If $v\in A$, a free vertex $w\in U_i$ stops being
free only if $i=j$ and $w\sim v$, that is, by (2), only if $w=v_j$. If $v$
lies in neither $A$ nor $U$, then by (1) nothing changes.

(c) \emph{The potential.} If a move makes a free vertex $y$ taken, every
alive $a\sim y$ stays alive with $t(a)$ one larger, so $\Phi$ grows by
$\Phi_y(c)$. If a move makes a free vertex $y$ dead, exactly the alive
$a\sim y$ stop being alive, so $\Phi$ falls by $\Phi_y(c)$. Otherwise
$\Phi$ is unchanged.

(d) \emph{Bob's strategy.} When some vertex is free, Bob takes a free
vertex $x$ with $\Phi_x(c)$ largest: if $x\in U_i$ he colors it $i$, which
is legal because $x$ is free. Otherwise he makes any legal move.

(e) \emph{The invariants.} For a position $c$ with Bob to move let
$I_B(c)$ say $8\Phi(c)>2^kF(c)$; for a position with Alice to move let
$I_A(c)$ say $8\Phi(c)>2^kF(c)$ and
$8(\Phi(c)-\Phi_y(c))>2^k(F(c)-1)$ for every free $y$.
At the start every vertex of $U$ is free, every $a\in A$ is alive with
$t(a)=0$, $\Phi=|A|$, $F=|U|$ and $\Phi_y\le D$, so $I_A$ holds by (4).
If $I_A(c)$ holds and Alice moves to $c'$, then by (b) and (c) either a
free $y$ became dead, with $\Phi(c')=\Phi(c)-\Phi_y(c)$ and
$F(c')=F(c)-1$, which is the second part of $I_A(c)$; or a free $y$ became
taken, and $\Phi$ grew while $F$ fell; or nothing changed. So $I_B(c')$
holds.

Now let $I_B(c)$ hold, let some vertex be free, and let Bob take $x$ as in
(d), reaching $c'$. Then $\Phi(c')=\Phi(c)+\Phi_x(c)$ and
$F(c')=F(c)-1$, which gives the first part of $I_A(c')$. Let $y$ be free in
$c'$; then $y\ne x$, $y$ was free in $c$, and $\Phi_y(c)\le\Phi_x(c)$. The
alive vertices of $c$ and $c'$ are the same, and $t$ changed only on those
adjacent to $x$. By (3) at most one alive $a$ is adjacent to both $x$ and
$y$; if there is one, $x$ and $y$ are two of $a_1,\dots,a_k$, both free in
$c$, so $t(a)\le k-2$ in $c$ and the term of $a$ in $\Phi_y$ grew by
$2^{t(a)}\le2^k/4$. Hence $8(\Phi_y(c')-\Phi_y(c))\le2\cdot2^k$ and
\[
8(\Phi(c')-\Phi_y(c'))\ge8\Phi(c)-2\cdot2^k>2^kF(c)-2\cdot2^k
=2^k(F(c')-1).
\]

(f) \emph{Bob wins.} Let Bob follow (d), and let $p$ be the first position
of the play with no free vertex; there is one, because the last position
has none (a free vertex can be colored). By (e) and induction along the
play, $I_A$ or $I_B$ holds at every position up to $p$, according to who
moves; the step for Bob is used only at positions with a free vertex. So
$\Phi(p)>0$ and some $a$ is alive at $p$. As no vertex is free,
$a_1,\dots,a_k$ are all taken, so $a$ sees all $k$ colors, and by (a) $a$
is uncolored. It can never be colored, and Bob wins.
\end{proof}

For $k\ge2$ this adapts the weighting argument behind Beck's criterion
\cite{Be}, \cite[Theorem~4.3]{Kr}: $U$ plays the role of the board and the
sets $\{a_1,\dots,a_k\}$ the winning sets of a $k$-uniform family in which
two sets share at most one element. New are the legality of Bob's moves,
tracked by the statuses free, taken and dead, and the overlap estimate
in (e). The term $8D$ pays for Alice's first move.

\begin{corollary}\label{cor:structure}
Let $K\ge2$ and let $G$ consist of pairwise disjoint vertex sets
$A,U_1,\dots,U_K$ and edges joining each $a\in A$ to exactly one vertex $a_i$
of each $U_i$, with no other edges (further isolated vertices are
allowed). If $G$ has no 4-cycle and $(2^K+8)|U|<8|A|$, where
$U=U_1\cup\dots\cup U_K$, then $\chi_g(G)>K$.
\end{corollary}

\begin{proof}
Fix $k\le K$ and apply Theorem~\ref{thm:potential} with the sets
$U_1,\dots,U_k$; the vertices of the other sets lie in neither $A$ nor the
union of the sets used. Conditions (1) and (2) are immediate, and (3)
holds because two vertices of $A$ with two common neighbors span a
4-cycle. For (4), let $w\in U_i$ and pick $j\ne i$. Two distinct neighbors
$a,a'$ of $w$ have $a_j\ne a'_j$, since otherwise $w$ and $a_j$ would be
two common neighbors. So $w$ has at most $|U_j|\le|U|$ neighbors, and (4)
holds with $D=|U|$ because $2^k|U_1\cup\dots\cup U_k|+8|U|\le(2^K+8)|U|$.
\end{proof}

\section{Extraction}\label{sec:extract}

\begin{lemma}\label{lem:mindeg}
A graph with at least one vertex and average degree at least $t>0$ has a
nonempty vertex set $S$ in which every vertex has at least $t/2$ neighbors
in $S$.
\end{lemma}

\begin{proof}
Write $e(S)$ for the number of edges inside $S$, and take $S$ nonempty of
least size with $2e(S)\ge t|S|$; the whole vertex set qualifies. If some
$v\in S$ had fewer than $t/2$ neighbors in $S$, then $S\setminus\{v\}$
would be nonempty, since $e(S)>0$, and
$2e(S\setminus\{v\})>2e(S)-t\ge t|S\setminus\{v\}|$, against the choice
of $S$.
\end{proof}

\begin{lemma}\label{lem:extract}
Let $K,M\ge0$ be integers and let $S$ be a nonempty set of vertices of a
graph such that each vertex in $S$ has at least $\delta=2(M+1)K(K+1)$
neighbors in $S$.
Then there are disjoint sets $U_1,\dots,U_K\subseteq S$ and a set
$A\subseteq S\setminus(U_1\cup\dots\cup U_K)$ such that every $a\in A$ has
a neighbor in each $U_i$ and $M|U_1\cup\dots\cup U_K|<|A|$.
\end{lemma}

\begin{proof}
For $K=0$ take $A=S$. Let $K\ge1$, $N=2(M+1)K$ and $n=|S|$, and consider
all $N^n$ labelings $\ell\colon S\to\{0,\dots,N-1\}$. For a labeling, let
$U_i(\ell)=\ell^{-1}(i-1)$ for $i\le K$, let $U(\ell)$ be their union, and
let $A(\ell)$ be the set of $x\in S$ with $\ell(x)\ge K$ that have a
neighbor in every $U_i(\ell)$. Each $x$ lies in $U(\ell)$ for $KN^{n-1}$
labelings, so $\sum_\ell|U(\ell)|=nKN^{n-1}$. For fixed $x$ and $i$, fix
$\delta$ neighbors of $x$ in $S$; the labelings in which none of them has
label $i-1$ number $(N-1)^\delta N^{n-\delta}$, and they include those in
which no neighbor of $x$ in $S$ has label $i-1$. So $x$ lies in $A(\ell)$ for at least
$N^n-KN^{n-1}-K(N-1)^\delta N^{n-\delta}$ labelings. Since $N\ge2$ and,
by the binomial theorem, $(N/(N-1))^N\ge1+N/(N-1)>2$, we have
$2(N-1)^N\le N^N$; as $\delta=N(K+1)$, this gives
$2^{K+1}(N-1)^\delta\le N^\delta$, and since $2^{K+1}>2K$ it follows that
$K(N-1)^\delta N^{n-\delta}<N^n/2$. Also $(M+1)KN^{n-1}=N^n/2$. Therefore
\[
\sum_\ell\bigl(|A(\ell)|-M|U(\ell)|\bigr)\ge
n\bigl(N^n-(M+1)KN^{n-1}-K(N-1)^\delta N^{n-\delta}\bigr)>0,
\]
and some labeling has $|A(\ell)|>M|U(\ell)|$.
\end{proof}


\section{The cited theorem}\label{sec:cited}

\begin{theorem}[{Kühn and Osthus \cite[Theorem~1]{KO}}]\label{thm:koorig}
``For every $k$ there exists $d=d(k)$ such that every graph of average
degree at least $d$ contains a subgraph of average degree at least $k$
whose girth is at least six.''
\end{theorem}

We use it in this form, the hypothesis \texttt{C4FreeSubgraphTheorem} of
the Lean proof (Section~\ref{sec:lean}).

\begin{theorem}\label{thm:ko}
For every integer $t\ge0$ there is an integer $d$ such that every graph
with at least one vertex and average degree at least $d$ has a subgraph
with at least one vertex, no 4-cycle and average degree at least $t$.
\end{theorem}

\begin{proof}[Derivation from Theorem~\ref{thm:koorig}]
For $t=0$ take $d=0$: a single vertex is such a subgraph. For $t\ge1$ take
$d=\lceil d(t)\rceil$ from Theorem~\ref{thm:koorig} with $k=t$. Since $d\ge d(t)$,
a graph of average degree at least $d$ has a subgraph of girth at least six
and average degree at least $t\ge1$. That
subgraph has an edge, hence a vertex, and its girth excludes 4-cycles.
\end{proof}
Appendix~\ref{app:bound} gives an explicit $d$.

\section{Proofs of the main results}\label{sec:proofs}

\begin{proof}[Proof of Theorem~\ref{thm:main}]
The conclusion for $K$ contains the conclusion for every smaller $K$, so
we may assume $K\ge2$. Let $M=2^K+1$, $\delta=2(M+1)K(K+1)$, $t=2\delta$,
and let $d$ be given by Theorem~\ref{thm:ko} for this $t$. Let $G$ have at
least one vertex and average degree at least $d$. Theorem~\ref{thm:ko}
gives a subgraph $G_0$ with no 4-cycle and average degree at least $t$;
Lemma~\ref{lem:mindeg} gives a set $S$ in which every vertex has at least
$\delta$ neighbors in $S$; and Lemma~\ref{lem:extract} gives $A$ and
$U_1,\dots,U_K$ with $M|U|<|A|$. Choose for each $a\in A$ one neighbor
$a_i\in U_i$ for each $i$, and let $H$ be the subgraph of $G_0$ on the
vertex set of $G_0$ whose edges are the pairs $aa_i$. Then $H$ has no
4-cycle, and $(2^K+8)|U|\le8M|U|<8|A|$, also when $U$ is empty, since then
$A$ is not. By Corollary~\ref{cor:structure}, $\chi_g(H)>K$.
\end{proof}

\begin{proof}[Proof of Corollary~\ref{cor:arboricity}]
Let $d=d(K)$ be as in Theorem~\ref{thm:main}, let $\mathcal C$ be a
monotone class whose members have game chromatic number at most $K$, and
let $G\in\mathcal C$. Every subgraph of $G$ lies in $\mathcal C$, so by
Theorem~\ref{thm:main} every subgraph of $G$ with at least one vertex has
average degree less than $d$. Hence every nonempty vertex set $S$ contains
a vertex with fewer than $d$ neighbors in $S$, since otherwise the subgraph
induced by $S$ would have average degree at least $d$. Removing such
vertices one at a time orders the vertices as $v_1,\dots,v_n$ so that each
$v_m$ has fewer than $d$ neighbors among $v_1,\dots,v_{m-1}$. Color the
edges from each $v_m$ to earlier vertices with distinct colors from
$1,\dots,d$. In each color class every vertex has at most one edge to an
earlier vertex, so the class has no cycle, as the latest vertex of a cycle
would have two. So $G$ has arboricity at most $d$.
\end{proof}

\section{Formal verification}\label{sec:lean}
Theorem~\ref{thm:main} and Corollary~\ref{cor:arboricity} are proved in
Lean~4 \cite{Lean} with the Mathlib library \cite{Mathlib}, from
definitions of the game given verbatim in Appendix~\ref{app:lean}. A position gives each vertex one of the $k$
colors, indexed by \texttt{Fin k}, that is $\{0,\dots,k-1\}$, or
\texttt{none}, meaning uncolored. A move gives one uncolored vertex a color
that none of its neighbors has. The inductive predicate
\texttt{AliceWins G k b c} says that Alice can force a win from the
position \texttt{c}, with Alice to move when \texttt{b} is \texttt{true}:
every vertex is colored; or it is Alice's turn and some move of hers leads
to such a position; or it is Bob's turn, he has a move, and all his moves
lead to such positions. \texttt{AliceWinsGame G k} is this predicate at
the empty position with Alice to move. On a finite graph,
\texttt{¬ AliceWinsGame G k} says exactly that Bob has a winning strategy
in the $k$-color game (Appendix~\ref{app:lean}). The cited theorem is the
proposition
\begin{verbatim}
def C4FreeSubgraphTheorem : Prop :=
  ∀ t : ℕ, ∃ d : ℕ, ∀ (V : Type) [Finite V] [Nonempty V] (G : SimpleGraph V),
    d * Nat.card V ≤ 2 * G.edgeSet.ncard →
    ∃ H : G.Subgraph, H.verts.Nonempty ∧ C4Free H.coe ∧
      t * H.verts.ncard ≤ 2 * H.coe.edgeSet.ncard
\end{verbatim}
where \texttt{C4Free H} says that there are no $a,b,c,d$ with $a\ne c$,
$b\ne d$ and $a\sim b\sim c\sim d\sim a$; this is Theorem~\ref{thm:ko},
with ``average degree at least $d$'' written as $d|V|\le2|E|$. The two
main results, in \texttt{Main.lean}, are
\begin{verbatim}
theorem GCG.exists_subgraph_bob_wins (hC4 : C4FreeSubgraphTheorem) (K : ℕ) :
    ∃ d : ℕ, ∀ (V : Type) [Finite V] [Nonempty V] (G : SimpleGraph V),
      d * Nat.card V ≤ 2 * G.edgeSet.ncard →
      ∃ H : G.Subgraph, ∀ k ≤ K, ¬ AliceWinsGame H.coe k

theorem GCG.bounded_arboricity (hC4 : C4FreeSubgraphTheorem) (K : ℕ) :
    ∃ d : ℕ, ∀ C : (V : Type) → SimpleGraph V → Prop,
      (∀ V G, C V G → Finite V) →
      (∀ V G, C V G → ∀ H : G.Subgraph, C H.verts H.coe) →
      (∀ V G, C V G → ∃ k ≤ K, AliceWinsGame G k) →
      ∀ V G, C V G → ArboricityLE G d
\end{verbatim}
The first is Theorem~\ref{thm:main}, where \texttt{H.coe} is the graph
formed by a Mathlib subgraph \texttt{H}. In the second, a class is a
predicate \texttt{C} whose three hypotheses say that its members are
finite, closed under subgraphs, and of game chromatic number at most $K$; \texttt{ArboricityLE G d} gives arboricity at most $d$
(Appendix~\ref{app:lean}). The build fails unless the axioms of
the two main theorems and of Theorem~\ref{thm:potential} are exactly
Lean's \texttt{propext}, \texttt{Classical.choice} and \texttt{Quot.sound}.
Theorem~\ref{thm:potential} and Lemma~\ref{lem:extract} are formalized as
written. Lemma~\ref{lem:mindeg} is checked for integer $t$, and
Corollary~\ref{cor:structure} under the stronger hypothesis
$(2^K+1)|U|<|A|$; these are the cases Section~\ref{sec:proofs} uses, and
Appendix~\ref{app:lean} notes the differences in detail and gives the
commands that reproduce the check.

\begin{appendices}
\section{The Lean definitions and correspondence}\label{app:lean}
A position is a function \texttt{c : V → Option (Fin k)}: \texttt{Fin k}
is the set $\{0,\dots,k-1\}$ indexing the $k$ colors, and the value
\texttt{none} means uncolored. The game is defined in
\texttt{GCG/Game.lean}:
\begin{verbatim}
def Move (G : SimpleGraph V) {k : ℕ} (c c' : V → Option (Fin k)) : Prop :=
  ∃ v i, c v = none ∧ (∀ w, G.Adj v w → c w ≠ some i) ∧ c' v = some i ∧
    ∀ u, u ≠ v → c' u = c u

inductive AliceWins (G : SimpleGraph V) (k : ℕ) :
    Bool → (V → Option (Fin k)) → Prop
  | full (b : Bool) (c : V → Option (Fin k)) :
      (∀ v, c v ≠ none) → AliceWins G k b c
  | alice (c c' : V → Option (Fin k)) :
      Move G c c' → AliceWins G k false c' → AliceWins G k true c
  | bob (c : V → Option (Fin k)) :
      (∃ c', Move G c c') → (∀ c', Move G c c' → AliceWins G k true c') →
      AliceWins G k false c

def AliceWinsGame (G : SimpleGraph V) (k : ℕ) : Prop :=
  AliceWins G k true (fun _ => none)
\end{verbatim}

\paragraph{Why \texttt{¬ AliceWinsGame G k} means that Bob wins.}
Each move colors a vertex, so on a finite graph every play ends after at
most $|V|$ moves. If \texttt{AliceWins} holds at a position, the
constructors used in its derivation give Alice a strategy that wins from
there: they never let Bob reach a position where it fails, and a play can
end only at a full coloring. Suppose it fails at a position. If the player
to move has no legal move, play ends there, and since \texttt{full} does
not apply, some vertex is uncolored: Bob has won. If it is Alice's turn,
then by \texttt{alice} it fails after each of her moves. If it is Bob's
turn and he has a move, then by \texttt{bob} it fails after some move of
his, and he plays that move. Following this, Bob keeps \texttt{AliceWins}
false until play ends, so he wins. Zhu \cite[pp.~1--2]{Zhu} awards Bob the
game as soon as an uncolored vertex has all $k$ colors among its
neighbors. Such a vertex stays uncolorable, so under the rules of
Section~1 play ends with it uncolored and Bob wins as well; conversely, a
play that ends with an uncolored vertex reached such a vertex. The two
conventions give the same winner.

\paragraph{Arboricity.} \texttt{ArboricityLE G d} says that some map from
all unordered pairs of vertices to $d$ colors makes every color class of
edges of $G$ a forest. Restricting the map to the edges partitions them
into at most $d$ forests, so $G$ has arboricity at most $d$; conversely,
for $d>0$ a partition into $d$ forests extends to such a map. For $d=0$
the predicate fails on every graph with a vertex, which does not affect an
existential upper bound.

\paragraph{Correspondence.}
Both main theorems take the cited theorem as the hypothesis \texttt{hC4},
and their proofs establish its hypotheses (finiteness, a vertex and the
density bound) for each graph to which they apply it.
\texttt{C4Free} is in \texttt{GCG/Game.lean} and \texttt{ArboricityLE} in
\texttt{GCG/Forest.lean}. Theorem~\ref{thm:potential} is
\texttt{GCG.bob\_wins\_of\_structure} in \texttt{GCG/Potential.lean}, in
which conditions (1)--(4) appear as hypotheses on a set \texttt{A}, a
family of classes and a map assigning to each $a\in A$ and $i$ its
neighbor $a_i$. Lemma~\ref{lem:mindeg} is
\texttt{GCG.exists\_minDegree\_subset} in \texttt{GCG/Extract.lean}, for a
positive integer threshold $t$, the case used in
Section~\ref{sec:proofs}; the written proof covers real $t>0$.
Lemma~\ref{lem:extract} is \texttt{GCG.extract} in the same file.
\texttt{GCG.bob\_wins\_pickGraph} in \texttt{Main.lean} is
Corollary~\ref{cor:structure} with the hypothesis $(2^K+8)|U|<8|A|$
replaced by the stronger $(2^K+1)|U|<|A|$, which is what
Lemma~\ref{lem:extract} provides in Section~\ref{sec:proofs}; the graph
whose edges are the pairs $aa_i$ is \texttt{GCG.pickGraph}. The printed
Corollary~\ref{cor:structure} follows from the checked
Theorem~\ref{thm:potential} by the written proof. The degeneracy step in
the proof of Corollary~\ref{cor:arboricity} is
\texttt{GCG.arboricityLE\_of\_degenerate} in \texttt{GCG/Forest.lean}.
The axiom checks are in \texttt{Axioms.lean}.

The formal proof of Theorem~\ref{thm:main} applies the cited theorem with
$t=4(2^{K+2}+2)(K+2)(K+3)$, which is the value of $t$ in
Section~\ref{sec:proofs} for $K+2$ in place of $K$; this covers $K<2$
without a separate case. Its subgraph $H$ keeps all the vertices of $G_0$,
as in the written proof.

To reproduce the check, run \texttt{lake exe cache get} and
\texttt{lake build} in the \texttt{lean} folder of the accompanying files,
which pins Lean v4.34.0 and a Mathlib commit; with the Mathlib cache the
build takes about a minute.

\section{Zhu's Questions 2 and 4}\label{app:related}
Kierstead and
Trotter \cite[Example~6.1]{KT} subdivide every edge of $K_{n,n}$ once;
their proof shows that these graphs $G_n$ have acyclic chromatic number at
most 3 and game chromatic number at most 4, and that their Go number tends
to infinity. By a theorem of Dinski and Zhu \cite[Theorem~2]{DZ}, every
subgraph of a graph with acyclic chromatic number at most $k$ has game
chromatic number at most $k(k+1)$. So the subgraphs of the graphs $G_n$
form a monotone class with game chromatic number at most 12 and unbounded
game coloring number, which answers Question~2 in the negative.

Dvořák \cite[Section~3.1.5]{Dv} showed that the conjecture of Dinski and
Zhu follows from Conjecture~3.1 of \cite{Dv} (Conjecture~1 of \cite{Dv08}, derived there from
its Corollary~4): ``There exists a function $f$ such that for each graph
$G$, if $\chi(G)\ge f(c)$ then the game chromatic number of
$\mathrm{sd}_1(G)$ is at least $c$'', where $\chi(G)$ is the chromatic
number, $c$ ranges over the positive integers and $\mathrm{sd}_1(G)$ is
obtained by subdividing every edge of $G$ once. Dvořák calls it ``wide
open''. The present argument does not reach these graphs. In a subgraph of
$\mathrm{sd}_1(G)$, every edge meets a subdivision vertex and these have
degree at most 2, so there are at most twice as many edges as subdivision
vertices; a subgraph with an edge also has a vertex of $G$, so its average
degree is less than 4, and Theorem~\ref{thm:main} says nothing about it.
\section{An explicit bound}\label{app:bound}
Montgomery, Pokrovskiy and Sudakov \cite[Theorem~1.1]{MPS} showed that
there is $c>0$ such that for every integer $t\ge1$, every graph of average
degree at least $2^{ct^2\log t}$ has a subgraph of average degree at least
$t$ and girth at least 6; so $d=\lceil2^{ct^2\log t}\rceil$ serves in
Theorem~\ref{thm:ko} for $t\ge1$. Section~\ref{sec:proofs} applies it with $t$ of order
$K^22^K$, so $d(K)$ is at most doubly exponential in $K$; we made no
attempt to optimize it.
\end{appendices}

\begin{thebibliography}{99}
\bibitem{Be} J. Beck, Van der Waerden and Ramsey type games,
\emph{Combinatorica} 1 (1981), 103--116,
\url{https://doi.org/10.1007/BF02579267}.
\bibitem{BGKZ} T. Bartnicki, J. Grytczuk, H. A. Kierstead and X. Zhu, The
map-coloring game, \emph{Amer. Math. Monthly} 114 (2007), 793--803.
\bibitem{DZ} T. Dinski and X. Zhu, A bound for the game chromatic number
of graphs, \emph{Discrete Math.} 196 (1999), 109--115. Theorem and
conjecture numbers as in the authors' preprint of October 1997,
\url{https://www.math.nsysu.edu.tw/~zhu/papers/game/bound.pdf}.
\bibitem{Dv} Z. Dvořák, \emph{Asymptotical Structure of Combinatorial
Objects}, doctoral thesis, Charles University, Prague, 2007,
\url{https://iuuk.mff.cuni.cz/~rakdver/disert.pdf}.
\bibitem{Dv08} Z. Dvořák, On forbidden subdivision characterizations of
graph classes, \emph{European J. Combin.} 29 (2008), 1321--1332,
\url{https://doi.org/10.1016/j.ejc.2007.05.008}. Numbering
as in the author's preprint, ITI Series 2006-324.
\bibitem{ES} P. Erdős and J. L. Selfridge, On a combinatorial game,
\emph{J. Combin. Theory Ser. A} 14 (1973), 298--301,
\url{https://doi.org/10.1016/0097-3165(73)90005-8}.
\bibitem{Ho} L. Hollom, \emph{On monotonicity in Maker--Breaker graph colouring
games}, \emph{Discrete Appl. Math.} 347 (2024), 131--138,
\url{https://doi.org/10.1016/j.dam.2024.01.007}.
\bibitem{KT} H. A. Kierstead and W. T. Trotter, Competitive colorings of
oriented graphs, \emph{Electron. J. Combin.} 8(2) (2001), \#R12.
\bibitem{Kr} M. Krivelevich, Positional games, in \emph{Proceedings of the
International Congress of Mathematicians, Seoul 2014}, vol.~4, Kyung Moon
Sa, 2014, 355--379. Theorem numbering as in arXiv:1404.2731v1.
\bibitem{KO} D. Kühn and D. Osthus, Every graph of sufficiently large
average degree contains a $C_4$-free subgraph of large average degree,
\emph{Combinatorica} 24 (2004), 155--162.
\bibitem{Lean} L. de Moura and S. Ullrich, The Lean 4 theorem prover and
programming language, in \emph{Automated Deduction -- CADE 28}, Lecture
Notes in Comput. Sci. 12699, Springer, 2021, 625--635.
\bibitem{Mathlib} The mathlib Community, The Lean mathematical library, in
\emph{Proceedings of the 9th ACM SIGPLAN International Conference on
Certified Programs and Proofs}, 2020, 367--381.
\bibitem{MPS} R. Montgomery, A. Pokrovskiy and B. Sudakov, $C_4$-free
subgraphs with large average degree, \emph{Israel J. Math.} 246 (2021),
55--71; arXiv:2004.03564v1.
\bibitem{TZ} Zs. Tuza and X. Zhu, \emph{Colouring games}, in \emph{Topics in
Chromatic Graph Theory} (L. W. Beineke and R. J. Wilson, eds.),
Encyclopedia Math. Appl. 156, Cambridge University Press, 2015, 304--326.
\bibitem{Zhu} X. Zhu, The game coloring number of planar graphs,
\emph{J. Combin. Theory Ser. B} 75 (1999), 245--258. Page and question
numbers as in the author's preprint,
\url{https://www.math.nsysu.edu.tw/~zhu/papers/game/planar.pdf}.
\end{thebibliography}
\end{document}
